% Notas internas para la presentación oral (no se incluye en main.tex)
\section{Notas para la presentación oral}
En la exposición conviene destacar ocho decisiones:
\begin{enumerate}
\item Por qué se escogió una cubierta institucional real con planos as-built.
\item Cómo la exigencia del ICE (UL 1741 SB / IEEE 1547-2018) y el NEC 690.4(B) obligaron a cambiar el inversor y el módulo.
\item Por qué el generador y el ATS llevaron la interconexión al lado normal del ATS, en lugar de usar la regla del 120\,\% en el TP, y por qué se agregó el interruptor secundario IS-TS01 de 400~A (NEC 240.4(F), 240.21).
\item Qué exige el ICE más allá de la ficha del inversor: informe ICE-LEE/ECA, categoría B con volt-var por defecto y SCADA si el abonado es binómico.
\item Cómo NFPA 1 y los lineamientos de viento del CFIA (Zona III, $V_b=115$~km/h, servicio y último por separado) definieron la huella del arreglo y la carga exigida a las abrazaderas.
\item Cómo se verificaron los strings con optimizadores: nivel módulo y reglas del fabricante para el string.
\item Por qué la estructura está caracterizada pero no certificada, y por qué el resultado es una factibilidad preliminar condicionada y no un cumplimiento definitivo.
\item Cómo se completó la información faltante del edificio con supuestos de diseño declarados (Cuadro 8) y por qué ninguno cambia la arquitectura del sistema.
\end{enumerate}
También se mencionará el uso de la herramienta de IA y la verificación de sus resultados contra las fuentes primarias.

